% =================================================================
\section{Introducción y alcance}
% =================================================================
CaribeShop es una tienda en línea de Cartagena que el año pasado tuvo tres problemas: se cayó
durante una promoción porque no aguantó la demanda, volvió a caerse porque dependía de una sola
máquina virtual, y una auditoría encontró la contraseña de la base de datos dentro del código y
el acceso remoto (SSH) abierto a Internet.

Este informe describe cómo diseñamos, desplegamos y probamos en Microsoft Azure una plataforma
que corrige esos tres problemas: es \textbf{segura}, \textbf{crece sola} cuando aumenta la demanda
y \textbf{sigue funcionando} cuando falla una de sus partes. Toda la infraestructura se creó con
código (Terraform) y cada requisito se comprobó con una prueba.

\begin{itemize}
  \item \textbf{Repositorio:} \url{https://github.com/AgEnt3P18/caribeshop-taller-integrador}
  \item \textbf{Herramientas:} Terraform, Azure CLI, GitHub Actions, Checkov (revisión de seguridad
        del código) y Locust (pruebas de carga).
  \item \textbf{Región:} East US, repartiendo los servidores en sus tres zonas de disponibilidad.
\end{itemize}

% =================================================================
\section{Arquitectura y decisiones de diseño}
% =================================================================
La Figura~\ref{fig:arq} resume el recorrido de una visita: entra por un balanceador de carga, que
la envía a uno de los servidores sanos; los servidores viven en una red privada, repartidos en
tres zonas, y consultan la contraseña en un Key Vault. Todo queda registrado para monitoreo.

\begin{figure}[H]
\centering
\begin{tikzpicture}[
  font=\small,
  caja/.style={draw=thead, line width=0.9pt, fill=white, rounded corners=4pt, minimum height=13mm,
               minimum width=31mm, align=center},
  serv/.style={fill=turq, text=white, rounded corners=3pt, minimum width=30mm, minimum height=9mm,
               font=\small\bfseries},
  flecha/.style={-{Stealth[length=2.6mm]}, line width=1.1pt, color=thead!80}
]
\node[caja] (u) {\textbf{Clientes}\\ desde Internet};
\node[caja, right=11mm of u] (lb) {\textbf{Balanceador}\\ \textbf{de carga}\\ \footnotesize revisa cada servidor\\ \footnotesize cada 5 segundos};
\node[serv, right=17mm of lb, yshift=12mm] (v1) {Servidor · zona 1};
\node[serv, below=3mm of v1] (v2) {Servidor · zona 2};
\node[serv, below=3mm of v2] (v3) {Servidor · zona 3};
\begin{scope}[on background layer]
  \node[draw=badred!70, dashed, line width=0.9pt, rounded corners=6pt, fill=badred!4,
        fit=(v1)(v3), inner xsep=5mm, inner ysep=4mm] (red) {};
\end{scope}
\node[font=\footnotesize\bfseries, color=badred!85, anchor=south] at ([yshift=0.5mm]red.north) {Red privada · solo entra tráfico web};
\node[font=\footnotesize\itshape, color=thead, anchor=north] at ([yshift=-1mm]red.south)
  {de 2 a 4 servidores, se agregan solos};
\node[caja, right=13mm of red.east, yshift=12mm] (kv) {\textbf{Key Vault}\\ \footnotesize guarda la contraseña};
\node[caja, right=13mm of red.east, yshift=-12mm] (mon) {\textbf{Monitoreo}\\ \footnotesize registros y alertas\\ \footnotesize por correo};
\draw[flecha] (u) -- (lb);
\draw[flecha] (lb) -- (red.west |- lb);
\draw[flecha] (red.east |- kv) -- (kv);
\draw[flecha] (red.east |- mon) -- (mon);
\node[fill=thead, text=white, rounded corners=4pt, inner xsep=4mm, inner ysep=2.2mm,
      font=\footnotesize, align=center, anchor=north] (gob)
  at ([yshift=-10mm]$(u.west |- red.south)!0.5!(kv.east |- red.south)$)
  {Todo se crea con \textbf{Terraform} \quad·\quad cada cambio lo revisa \textbf{GitHub Actions + Checkov}\\
   \textbf{Azure Policy} bloquea lo prohibido \quad·\quad el estado de Terraform vive protegido en Azure};
\end{tikzpicture}
\caption{Arquitectura desplegada en Azure.}
\label{fig:arq}
\end{figure}

\begin{table}[H]
\centering\small\tablaestilo
\caption{Principales decisiones de diseño.}
\label{tab:decisiones}
\begin{tabularx}{\textwidth}{|>{\raggedright\arraybackslash\bfseries}p{4.6cm}|Y|}
\hline
\encab \hd{Qué elegimos} & \hd{Por qué} \\ \hline
Servidores en las 3 zonas de East US & Si una zona falla, las otras siguen atendiendo. Elegimos el tamaño D2s\_v4 porque era el que tenía cupo disponible en las tres zonas. \\ \hline
Balanceador de carga Standard & Revisa la salud de cada servidor y deja de enviarle tráfico si falla. Para producción se recomienda además un firewall de aplicaciones web (WAF). \\ \hline
Aplicación sin librerías externas & Así los servidores no necesitan salir a Internet, y pudimos \textbf{bloquear toda la salida}. Usa un proceso por núcleo para aprovechar toda la CPU. \\ \hline
Contraseña en Key Vault & La aplicación la lee con su propia identidad y \textbf{solo tiene permiso sobre ese único secreto}. \\ \hline
Bloqueo de salida con la red (NSG) & Solo se permite hablar con Key Vault, Monitor y Storage. Es gratis, frente a $\approx$\,USD 900/mes de un firewall dedicado. \\ \hline
Reparación automática & Si un servidor queda enfermo más de 10 minutos, Azure lo reemplaza solo. \\ \hline
Dos políticas de Azure (Deny) & Prohíben servidores con IP pública y recursos fuera de East US, sin importar desde dónde se creen. \\ \hline
\end{tabularx}
\end{table}

% =================================================================
\section{Implementación paso a paso}
% =================================================================
\begin{enumerate}
  \item \textbf{Acceso a Azure.} Iniciamos sesión con la CLI. El método por código de dispositivo
        estaba bloqueado por la seguridad predeterminada de la cuenta, así que usamos el inicio
        de sesión por navegador.
  \item \textbf{Región y tamaño.} Revisamos las cuotas de la suscripción (10 núcleos por región)
        y qué tamaños estaban disponibles en las tres zonas.
  \item \textbf{Estado de Terraform.} Creamos primero una cuenta de almacenamiento protegida, sin
        llaves de acceso, para guardar el estado de Terraform fuera de Git.
  \item \textbf{Despliegue.} Con \texttt{terraform plan} y \texttt{apply} se crearon los 33 recursos
        de una sola vez.
  \item \textbf{Configuración automática.} Cada servidor se instala solo al arrancar
        (\emph{cloud-init}); cuando cambiamos la aplicación bastó con volver a aplicar y reiniciar
        las instancias, sin entrar a ninguna máquina.
  \item \textbf{Pruebas y evidencias.} Ejecutamos las pruebas P1, P2 y P3 y guardamos los resultados
        en la carpeta \texttt{evidencias/} del repositorio.
\end{enumerate}

\textbf{Por qué el autoscale usa umbrales distintos para subir y bajar.} Agrega un servidor cuando
la CPU pasa del 70\,\% durante 5 minutos y quita uno cuando baja del 30\,\% durante 10 minutos.
Entre ambos valores no hace nada, lo que evita que agregue y quite servidores sin parar. Bajar es
más lento a propósito: quitar capacidad antes de tiempo es peor que mantenerla unos minutos más.
Tras cada acción espera 5 minutos (al subir) o 10 (al bajar), el tiempo que tarda una instancia
nueva en arrancar y estabilizarse.

% =================================================================
\section{Trazabilidad de requisitos}
% =================================================================
\begin{table}[H]
\centering\footnotesize\tablaestilo
\caption{Cada requisito, dónde se implementa y cómo se comprobó.}
\label{tab:traz}
\begin{tabularx}{\textwidth}{|>{\raggedright\arraybackslash}p{4.2cm}|>{\raggedright\arraybackslash}p{3.7cm}|Y|}
\hline
\encab \hd{Requisito} & \hd{Dónde se implementa} & \hd{Evidencia} \\ \hline
\grupo{3}{Seguridad} \\ \hline
S1 · Todo en Terraform & \texttt{providers.tf}, \texttt{bootstrap/} & Estado remoto; Git sin estado ni credenciales \\ \hline
S2 · Red mínima & \texttt{network.tf} (NSG) & Servidores sin IP pública; SSH inalcanzable \\ \hline
S3 · Key Vault & \texttt{keyvault.tf} & \texttt{/secret} muestra solo longitud y huella \\ \hline
S4 · Mínimo privilegio & \texttt{keyvault.tf} & Un único rol, limitado al secreto \\ \hline
S5 · Azure Policy & \texttt{policy.tf} & Rechazo \texttt{RequestDisallowedByPolicy} (Lst.~\ref{lst:policy}) \\ \hline
S6 · Auditoría & \texttt{monitoring.tf} & Consulta KQL de lecturas (Lst.~\ref{lst:kql}) \\ \hline
\grupo{3}{Escalamiento y alta disponibilidad} \\ \hline
E1 · Tres zonas & \texttt{compute.tf} & Instancias en zonas 1, 2 y 3 \\ \hline
E2 · Sonda de salud & \texttt{loadbalancer.tf} & Métrica de salud (Fig.~\ref{fig:p2logica}) \\ \hline
E3 · Autoscale & \texttt{compute.tf} & Historial 2\,$\rightarrow$\,4\,$\rightarrow$\,2 (Fig.~\ref{fig:stair}) \\ \hline
E4 · Reparación automática & \texttt{compute.tf} & Instancia reemplazada sola (P2.1) \\ \hline
E5 · Alertas por correo & \texttt{monitoring.tf} & Historial de alertas (Tabla~\ref{tab:alertas}) \\ \hline
\grupo{3}{DevSecOps} \\ \hline
D1 · Pipeline con Checkov & \texttt{.github/workflows/} & PR \#1 bloqueado y PR \#2 aprobado \\ \hline
D2 · Riesgos justificados & comentarios en cada recurso & Tabla~\ref{tab:riesgos} \\ \hline
D3 · Instalación automática & \texttt{cloud-init} + \texttt{app/} & Servidores listos sin pasos manuales \\ \hline
\end{tabularx}
\end{table}
